\section{Discussion}
\label{sec:discussion}

\subsection{Key Findings}

\textbf{The architecture-family adversarial fingerprint is real at base model scale, but
statistically underpowered at $N=9$.} $\eta^2=0.293$ ($f^2 \approx 0.41$, large) is consistent
across 83\% of attack categories. This is not a borderline effect: our estimate substantially
exceeds the $f^2>0.35$ large-effect boundary. The finding is consistent with
\citet{Neerudu2023Robustness}, who observed directional architecture-family differences in GLUE
robustness, and provides the first formal effect-size quantification of this phenomenon.

\textbf{adv\_rte (NLI paraphrase) produces the strongest family differentiation.}
$\eta^2=0.592$ ($p=0.075$, near-significant at $N=9$) suggests that adversarial
entailment/paraphrase detection is where architecture family effects concentrate. The
$\Delta^*$-vector pipeline itself --- 7 validated, reusable modules --- is a methodological
contribution independent of the statistical outcome.

\subsection{Limitations}

\paragraph{Statistical underpowering (primary).}
$N=9 \to \sim$40\% power at $\eta^2=0.29$. The $p=0.147$ result is expected from the power
analysis; the $\eta^2=0.293$ effect size stands as a practically meaningful finding. The h-e1-v2
study with $N=15$ directly addresses this with 80\% power. \emph{Suggested framing:} Our PoC
establishes the effect magnitude and provides power analysis guidance ($N \geq 15$), advancing
the design of future definitive studies.

\paragraph{LOMO classification N-degeneracy.}
LOMO=0.333 at exact chance. At $N=3$/family, cosine $k=1$ nearest-neighbor classification is
geometrically near-degenerate; at-chance LOMO is mathematically expected. The LOMO test at
$N=3$/family is not a valid test of the classification hypothesis. Running LOMO at $N=15$
(h-e1-v2) directly tests whether accuracy rises above chance as the degeneracy is resolved.

\paragraph{Mechanism untested.}
The attention-topology mediation hypothesis (h-m1--h-m4: $\Delta C$ extraction and mediation
test) was not executed, as it was gated on h-e1 gate satisfaction. The descriptive finding
($\eta^2=0.293$) and methodological contribution (pipeline) are independent of mechanism
confirmation.

\paragraph{enc\_dec ANLI-R3 coverage gap.}
T5/BART were fine-tuned only on SST-2, not MNLI. Their ANLI-R3 clean accuracy is near chance,
producing $\Delta^* \approx 0$ and $\eta^2=0.0$ --- a task-coverage gap, not a genuine null.
Corrective step (MNLI fine-tuning for T5/BART) is included in h-e1-v2.

\paragraph{CheckList partition skipped.}
CheckList behavioral testing was not evaluated due to package availability. Cross-partition
replication is partial; full replication requires CheckList inclusion in h-e1-v2.

\subsection{Broader Impact}

\textbf{Positive:} The $\Delta^*$-vector pipeline provides a practical tool for robustness
auditing at the architecture-family level, enabling principled model selection for
adversarial-sensitive deployments. The validated seven-module codebase enables community
replication and extension to new architectures, scales, and attack types.

\textbf{Potential concerns:} Adversarial vulnerability characterization is dual-use. We note
that (a) our characterization operates at the family level, not the individual-model level, and
(b) all attack methods we evaluate are publicly available benchmarks. We do not introduce new
attacks or disclose novel vulnerability types beyond the public benchmark literature.

\textbf{Scale effects:} Results are specific to base-scale models ($\sim$110--350M parameters).
Whether architecture-family fingerprints persist at 7B+ scale is an open empirical question.
RLHF and instruction-tuning are known to change robustness profiles; our results do not extend
to RLHF-tuned models.
